
%!TEX root = ../main.tex

\chapter{Grundlagen}
\label{chap:grundlagen}

\section{ATLAS aus Anwender-Sicht}
In diesem Abschnitt wird \ac{ATLAS} aus Sicht der Anwender beschrieben. Technische Details werden hier bewusst nur kurz angesprochen und in \autoref{chap:ist-analyse-anforderungen} aus Entwickler-Sicht genauer erklärt.

\ac{ATLAS} ist eine Webanwendung, die auf einem Angular-Frontend und einem .NET-basierten Backend aufbaut. Aus Anwender-Sicht dient \ac{ATLAS} als Plattform zur Dokumentation von Laser-Applikationen und Versuchen und unterstützt damit direkt die Anwendungsentwicklung sowie die Kundenberatung im \ac{LAC}. Eine saubere und vollständige Datenbasis ist dabei entscheidend, da fehlerhafte Einträge die Entscheidung bei Experimenten und Beratung der Kunden beeinträchtigen können.

Aktuell ist das System auf Experimente im Bereich Laserschweißen ausgerichtet. Berechtigte Nutzer können eigene Experimente erstellen und alle relevanten Versuchsinformationen strukturiert erfassen, zum Beispiel die Schweißkonfiguration, verwendete Materialien, Maschine, Laser sowie allgemeine Parameter wie Laserleistung und Schweißgeschwindigkeit. Ein Experiment besteht aus mehreren Versuchen, bis ein passendes Ergebnis erreicht wird.

Zusätzlich erzeugt das System automatisiert einen Report als PowerPoint-Präsentation, in den definierte Informationen und Bilder übernommen werden. Dieser Report kann für die Kundenkommunikation genutzt oder durch den Sales-Bereich weiterverwendet werden.

Welche Informationen ein Nutzer sehen oder bearbeiten darf, hängt von der zugewiesenen Rolle ab.

\subsection{Ablauf beim Anlegen eines Experiments}
Das Anlegen und Bearbeiten eines Experiments folgt in \ac{ATLAS} einer klaren Schrittfolge:

\begin{enumerate}
	\item \textbf{Experiment -- allgemeine Informationen:} Erfassung von Stammdaten wie Projekt-ID, Projektname, Confidentiality Level, zuständiger Engineer und weiteren Basisinformationen.
	\item \textbf{Experiment Task:} Definition des technischen Versuchsrahmens, zum Beispiel Type of Joint (z. B. Butt Joint, T Joint), Type of Weld (z. B. I-Weld, Fillet Weld), Weld Type (z. B. Partial Penetration, Full Penetration) sowie Materialangaben.
	\item \textbf{Setup:} Auswahl und Konfiguration von Maschine, Laser, Optik und weiteren spezifischen Prozessparametern.
	\item \textbf{Parameters:} Übersicht der einzelnen Versuche mit den jeweils gewählten Parametern, die der Nutzer in seiner Ansicht sehen und vergleichen möchte.
	\item \textbf{Result:} Aufbereitung der Ergebnisse einschließlich automatisiertem PowerPoint-Report mit den relevanten Informationen und eingebundenen Bildern.
\end{enumerate}

Diese Struktur unterstützt eine nachvollziehbare Dokumentation von Experimenten und erleichtert die spätere Auswertung.

\subsection{Rollen und Berechtigungen}
In \ac{ATLAS} existieren mehrere Rollen mit unterschiedlichen Funktions- und Sichtrechten. Die Rechte werden rollenspezifisch vergeben und in einzelnen Bereichen zusätzlich durch den Kontext (z. B. Owner/Non-Owner) eingeschränkt.

\textbf{ADMIN:} Die Rolle verfügt über vollständige Rechte innerhalb der Anwendung und kann alle relevanten Funktionen nutzen.

\textbf{SALES:} Die Rolle hat einen stark eingeschränkten Zugriff. Sie kann nur ausgewählte Informationen zu Experimenten sehen, keine Experimente erstellen und keine Detailinformationen innerhalb eines Experiments einsehen. Zusätzlich ist der Zugriff auf Berichte auf den Sales-Report begrenzt. Für Suchfunktionen werden nur die jeweils freigegebenen Informationen angezeigt.

Weitere Rollen wie App-Engineer, R\&D Engineer und External Engineer besitzen ebenfalls eigene, funktionsbezogene Rechteprofile.

 \todo{Bild von Test Atlas UI einfügen}



\section{Stand der Technik}
Für den Aufbau einer Testinfrastruktur in \ac{ATLAS} sind bestimmte Technologien, Werkzeuge und Konzepte grundlegend. Dieser Abschnitt gibt einen Überblick über die wichtigsten davon und erklärt jeweils, welche Rolle sie im Testprozess spielen. Die genaue Einordnung in die Systemarchitektur von \ac{ATLAS} folgt in \autoref{chap:ist-analyse-anforderungen}.

\subsection{Softwaretests im Überblick}
Automatisierte Softwaretests lassen sich in drei Hauptebenen einteilen: Unit-Tests prüfen einzelne Komponenten isoliert, Integrationstests überprüfen das Zusammenspiel mehrerer Teile, und End-to-End-Tests simulieren vollständige Nutzerabläufe. Die Testpyramide beschreibt dabei, wie diese Ebenen sinnvoll gewichtet werden sollten: viele schnelle Unit-Tests als Basis, weniger Integrationstests und eine kleine Zahl an End-to-End-Tests an der Spitze \cite{Cohn2009}. Grundlage für die Strukturierung von Testprozessen bildet zudem \cite{IEEE29119_2_2021}. Wichtig für Unit-Tests ist außerdem die klare Trennung von Abhängigkeiten über sogenannte Test Doubles \cite{Fowler2007}.

\subsection{Azure DevOps und Azure Blob Storage}
Azure DevOps ist eine Plattform von Microsoft für die Verwaltung von Softwareprojekten. Im Kontext dieser Arbeit ist vor allem Azure Pipelines relevant, das Build-, Test- und Deployment-Prozesse automatisiert \cite{MicrosoftAzurePipelines2026,MicrosoftAzurePipelinesArchitecture2026}. Azure Blob Storage dient zur Speicherung von Dateien und Artefakten in der Cloud. Für lokale Tests ohne echte Cloud-Anbindung bietet Microsoft mit Azurite einen Emulator, der das Speicherverhalten nachbildet \cite{MicrosoftAzurite2026}.

\subsection{Visual Studio}
Visual Studio ist die primäre Entwicklungsumgebung für das Backend von \ac{ATLAS}. Neben der normalen Entwicklung bietet es integrierte Werkzeuge für Tests: Der Test Explorer ermöglicht das Ausführen und Auswerten von Unit-Tests direkt in der IDE, Code-Coverage zeigt, welche Codebereiche durch Tests abgedeckt sind \cite{MicrosoftTestExplorerBasics2026,MicrosoftTestingToolsOverview2026}.

\subsection{Angular}
Angular ist ein TypeScript-basiertes Framework von Google für die Entwicklung von Webanwendungen und bildet das Frontend von \ac{ATLAS}. Für Tests stellt Angular das TestBed bereit, das eine isolierte Testumgebung für Komponenten und Services schafft. Als Test-Framework wird Jasmine eingesetzt, als Test-Runner Karma \cite{AngularTestingKarma2026,JasmineDocs2026,KarmaDocs2026}.

\subsection{Docker}
Docker dient zur Containerisierung von Anwendungen. Ein Container bündelt eine Anwendung mit allen ihren Abhängigkeiten und läuft dadurch auf jedem System gleich. Mit Docker Compose lassen sich mehrere Container gemeinsam verwalten, was das Aufsetzen von Testumgebungen vereinfacht \cite{DockerDocs2026,DockerComposeSpec2020}. Containeransätze eignen sich grundsätzlich gut für reproduzierbare und skalierbare Testinfrastrukturen \cite{Lemos2026}.

\subsection{.NET Core}
.NET Core ist ein quelloffenes Framework von Microsoft für die plattformübergreifende Entwicklung von Anwendungen mit C\#. Das Backend von \ac{ATLAS} basiert auf ASP.NET Core. Für Tests bietet das Framework unter anderem die \texttt{WebApplicationFactory}, mit der sich die Anwendung in einem In-Memory-Testserver starten lässt, ohne einen echten Server zu benötigen \cite{MicrosoftTestingASPNETCore2026,MicrosoftTestASPNETCoreMVC2026}. Empfehlungen zur Teststruktur und Benennung sind in den offiziellen Best Practices beschrieben \cite{MicrosoftUnitTestingBestPractices2026}.

\subsection{Frameworks}
Neben den genannten Technologien kommen im Projekt konkrete Test-Frameworks zum Einsatz, die im Folgenden kurz vorgestellt werden.

\subsubsection{xUnit}
xUnit ist ein weit verbreitetes Open-Source-Framework für Unit-Tests in .NET und wird von Microsoft in der offiziellen Dokumentation empfohlen. Im Vergleich zu Alternativen wie NUnit oder MSTest zeichnet sich xUnit durch eine modernere Architektur, bessere Unterstützung für Dependency Injection und parallele Testausführung aus \cite{xUnitDocs2026,MicrosoftUnitTestingBestPractices2026}.

\subsubsection{Playwright}
Playwright ist ein Framework von Microsoft für automatisierte End-to-End-Tests im Browser. Es unterstützt Chromium, Firefox und WebKit über eine gemeinsame API und wartet automatisch auf Elemente, was Tests stabiler macht als ältere Ansätze \cite{PlaywrightDocs2026}.

\subsection{VM}
Eine virtuelle Maschine (VM) ist ein softwarebasiertes Abbild eines Computersystems, das auf echter Hardware läuft. Sie verhält sich wie ein eigenständiger Rechner mit eigenem Betriebssystem und isolierter Umgebung. Für Testinfrastrukturen sind VMs interessant, weil sie einen definierten, jederzeit wiederherstellbaren Ausgangszustand bieten. Mit Tools wie Vagrant lassen sich VMs automatisiert provisionieren \cite{VagrantDocs2026,Paleologu2008}.



\section{Verknüpfung zu möglichen Vorlesungsinhalten}
softwareengineering (vorlseung welche testarten es gibt), datenbanken vorlesung (Postgresql, weil dort daten gespeichert werden)